% Pile lithium-iode d'un stimulateur cardiaque, vue en coupe (exercice 12, chapitre 8)
\begin{tikzpicture}
\filldraw[fill=otRule!50,draw=otInk!75,line width=0.7pt,rounded corners=2mm] (0,0) rectangle (5.2,2.6);
\fill[violet!45!black!65,rounded corners=1.2mm] (0.18,0.18) rectangle (5.02,2.42);
\fill[pcJaune] (1.5,0.55) rectangle (3.7,2.05);
\filldraw[fill=otInk!14,draw=otInk!60,line width=0.4pt] (1.62,0.67) rectangle (3.58,1.93);
\fill[otInk!70] (2.54,1.93) rectangle (2.66,3.25);
\filldraw[fill=white,draw=otInk!70,line width=0.4pt] (2.38,2.42) rectangle (2.82,2.8);
\draw[pcFil] (4.6,2.6) -- (4.6,3.25);
\node[font=\otSemi\fontsize{10}{11}\selectfont,text=otInk] at (2.6,3.5) {$-$};
\node[font=\otSemi\fontsize{10}{11}\selectfont,text=otInk] at (4.6,3.5) {$+$};
\pcNote[west]{(3.4,1.3)}{(6.0,2.45)}{anode : lithium}
\pcNote[west]{(3.64,0.8)}{(6.0,1.75)}{électrolyte solide : iodure de lithium}
\pcNote[west]{(4.5,0.5)}{(6.0,1.05)}{cathode : diiode et polymère conducteur}
\pcNote[west]{(5.2,0.3)}{(6.0,0.35)}{boîtier métallique, relié à la borne $+$}
\end{tikzpicture}
\par\smallskip\legende{Coupe schématique d'une pile lithium-iode : le lithium, au centre, est entouré par le mélange de diiode.}
